\section{Changing a raw lattice on the original base dictionary}

% BEGIN V74M1
\paragraph{V74M1. Simultaneous roots for a rational column move.}
\label{claim:V74M1}
Let $M=l.M$, $\mathfrak b=l.Mb$,
$\mathfrak b\bar{\mathfrak b}=(M)$ and
$\mathfrak b+\bar{\mathfrak b}=\mathcal O_E$. Let $L$ be a marked
target period lattice with marking $m:D_w.\mathrm{La}\to L$.
Let $d$ be a rational unitary column change integral with integral
inverse over $M$. Let $N,N'$ be the two positive allowed short
levels, and let $\pi$ be the induced bijection on
$\mathcal Q=(D_w.\mathrm{La}/\mathfrak bD_w.\mathrm{La})^n$.
In applying the complex short-value formula one may choose the
original roots of both markings so that
\[
                  Y'(\pi q)=dY(q)\qquad(q\in\mathcal Q)
\]
as rational matrices, including when $N\ne N'$.

\begin{proof}
Represent the entries of $q$ by
$a\in(\bar{\mathfrak b}D_w.\mathrm{La})^n$, possible by
comaximality. At $P\mid\mathfrak b$ prescribe the coset
$Y(q)\in M^{-1}m(a)+L_P^n$.
At each $P\nmid\mathfrak b$ require
\[
       Y(q)\in 2N L_P^n\cap d^{-1}(2N'L_P^n).
\]
These local conditions equal the integral lattice away from
finitely many places. They admit a rational solution. Indeed a
single common denominator puts all cosets in one fractional
lattice, and each condition then specifies a congruence in that
lattice modulo a power of a prime. The Chinese remainder theorem
on its finite quotient solves the finite list simultaneously.

At $P\nmid\mathfrak b$ the vector $m(a)/M$ is integral: at
$P\mid\bar{\mathfrak b}$ its numerator is divisible by the full
$P$-part of $M$. Thus $Y(q)-m(a)/M\in L^n$. All places outside
$\mathfrak b$ already satisfy $Y(q)\in 2N L_P^n$;
increasing $j$ at the finitely many remaining places gives
$\mathfrak b^jY(q)\subset 2N L^n$. These are exactly the two
root conditions. Choose a half representative $a'$ of $\pi q$.
At each $P\mid\mathfrak b$, $a'-da\in M(D_w.\mathrm{La})_P^n$
since the localized half modulus there is $(M)$. Linearity of $m$
after localization gives $dY(q)-m(a')/M\in L_P^n$. Away from
this half the two terms are integral, and $dY(q)\in 2N'L_P^n$
by its construction. Increasing $j$ again supplies the second
root condition with $N'$. Invertibility of the residue action
makes $Y'$ a definition for every $q'$. The point formula permits
every admissible system of roots, so using this pair introduces
no new invariance assumption.
\end{proof}
% END V74M1

% BEGIN V74M2
\paragraph{V74M2. One unit for the projected vacua seen by the base dictionary.}
\label{claim:V74M2}
Use the genuine two-group Weil action of the pinned base tuple.
Let $S=\mathcal O_E^n$, $L_*$ be the target lattice of the base
vacuum $v_*$, and $v_k=\omega(1,k)v_*$ for a finite target mover.
Write $L_k=kL_*$ at the finite completions. For any of the self-dual
source lattices $\Gamma$ under consideration define $p_{\Gamma,k}$
to be the orthogonal projection onto the prescribed finite vacuum
character of $\widehat\Gamma\otimes\widehat L_k$; put
$u_{\Gamma,k}=p_{\Gamma,k}v_k/\|p_{\Gamma,k}v_k\|$.
For a rational unitary column move $d$, one scalar of modulus one,
independent of the target mover, intertwines the normalized vacua
for $\Gamma$ and $d\Gamma$ on all these points.

\begin{proof}
The projector is the probability Haar average of the finite
Heisenberg operators multiplied by the prescribed semicharacter.
The latter is the finite additive character of one half the
alternating pairing of $y$ with $w_0 y$, with the pinned sign.
Its restriction to rational columns is the factor
$\exp(\pi i\sum_b E_F(y_b,w_0 y_b))$ in the classical projector;
this follows from the product formula for the additive character.
The self-dual integral lattice makes the corrected operators a
group on that lattice. On a smooth vector the integral is a finite
orbit average. Density of the rational tensor lattice in its finite
completion identifies it with the original finite sum, with its
original denominator. In particular this comparison commutes with
reading the adelic vacuum in the classical model of P74I5.

For $k=1$ let $A=\widehat\Gamma\otimes\widehat L_*$ and
$B=\widehat S\otimes\widehat L_*$. Representatives of
$A/(A\cap B)$ give the summands of $p_{\Gamma,1}v_*$. They
have norm one. If two representatives differ outside $B$, an
element of $B$ separates their Heisenberg characters by self-duality;
unitarity then makes the two vectors orthogonal. Representatives
differing by $A\cap B$ give the same corrected summand. Hence the
squared norm of the average is $[A:A\cap B]^{-1}$, which is
positive. Writing $r=\operatorname{rank} L_*$, this index equals
$[\Gamma:\Gamma\cap S]^r$: locally $L_*$ is free of rank $r$
and the local indices multiply.

At the base, the simultaneous eigenspace for this finite vacuum
and the archimedean Gaussian equations is a line. In the
Schr\"odinger realization the two halves of the self-dual
Heisenberg lattice impose the support and translation equations:
the diagonal half restricts the function to one dual coset and
the translation half determines all its values there from one
value. One may choose integral symplectic coordinates at each
finite prime; self-duality supplies the dual pair and the fixed
semicharacter supplies the coset and the translation factors.
At infinity the Gaussian equations determine a scalar multiple
of the normalized Gaussian. Equivalently this is the
one-dimensional theta space for the principal vacuum polarization
used in P74I5. This argument applies to $A$ and to its image by $d$.

Use the rational source element $a=(d^{-1})^{\mathsf T}$; the
unitary relation gives $a=\bar d$. The native action is right
multiplication by $a^{-1}$ on maps from $V$ to $W$, hence it is
left multiplication by $d$ on their transposed matrix of rows.
It carries $A$ to the lattice for $d\Gamma$, preserves the
semicharacter by $E$-linearity and isometry, and preserves the
Gaussian equations since the source unitary group is compact at
infinity. Thus the unit vectors $\omega(a,1)u_{\Gamma,1}$ and
$u_{d\Gamma,1}$ are proportional. Define the unit $\tau$ by
$u_{d\Gamma,1}=\tau\,\omega(a,1)u_{\Gamma,1}$.

The finite target action also preserves the very same semicharacter.
Conjugating the probability average by $\omega(1,k)$ therefore
carries $p_{\Gamma,1}$ to $p_{\Gamma,k}$, including its
normalization. It follows that
$u_{\Gamma,k}=\omega(1,k)u_{\Gamma,1}$ for both source
lattices. The source and target actions commute because the
original action is a representation of their direct product.
Consequently
\[
          u_{d\Gamma,k}=\tau\,\omega(a,1)u_{\Gamma,k}
\]
with the one unit just chosen. The scalar comparing the classical
standard vacuum with $v_k$ in P74I5 multiplies both classical
unit projections by the same phase. It cancels, even if it
varies between components.
\end{proof}
% END V74M2

% BEGIN V74M3
\paragraph{V74M3. The raw short-value change at the required dictionary points.}
\label{claim:V74M3}
For the original base dictionary and the rational moves of S74B3
and S74B4, let $U,U'$ be the normalization-one short readings with
corresponding Fourier weights and with $P'$ the column pushforward
of $P$. On all those dictionary points the same $\tau$ of V74M2
satisfies
\[
               U'\sqrt{j}=\tau\sqrt{j'}\,U,\qquad
 j=[\Gamma:\Gamma\cap S]^r,\qquad
 j'=[d\Gamma:d\Gamma\cap S]^r .
\]
Here the levels $N,N'$ may differ, the two geometric raw sections
may have different Fourier weights, and zero short values are
allowed.

\begin{proof}
On a marked analytic fiber take the same period frame $x$ and
the same product-to-classical generator constant $\kappa_x$ for
both sides as in P74I5. In the next display $u_{\Gamma,k}$ denotes
the corresponding classical normalized projection. To pass from the
adelic reading to this convention, divide by the factor comparing
$v_k$ to the classical unit standard vacuum of this chart. Projection
is linear and its norm on each side changes by the same factor
$1/\sqrt j$, so this comparison factor is common to both source
lattices. Let $Z_N$ denote the finite isotropic
subgroup in the short-value sum, and let $m_N$ be its normalization
index $\mathrm{mIdx}(\Gamma,N,r)$. The short reading of T73C1
in this marking is
\[
 U=\frac{\kappa_x\sqrt{j}}{|\mathcal Q|}
      \sum_{q\in\mathcal Q} w(q)\,
          \mathrm{thP}_x(T_{Y(q)}u_{\Gamma,k},P).
\]
For the index factor, the original geometric sum uses denominator
$m_N|\mathcal Q|$. Reading the translated theta factors and summing
over $Z_N$ gives its cardinality times the vacuum projection: each
term is the weighted Heisenberg translate for its isotropic point
and all fibers of the finite averaging map have equal size. In
integral tensor coordinates $|Z_N|=m_N j$, by the index tower
through $\Gamma\cap S$ and local freeness of the rank-$r$ target
lattice. The norm of that projection is $1/\sqrt j$ by V74M2.
The three factors give the displayed coefficient. The connection in
the short sum contributes exactly the classical harmonic derivative
$\mathrm{thP}_x$: in normal coordinates Taylor expansion of the
translated product applies the indicated directional derivatives;
terms involving derivatives of its Gaussian gauge contract two
letters and vanish on the harmonic polynomial. This is the same
short-value reading used for a class in P74I6--P74I7, here applied
before any assumption on the split support of a representative.

Use V74M1 for $Y,Y'$. At infinity let $d z$ denote the column
change on each Lie-type block, using that block's embedding of $E$.
In the adelic theta model a rational pair acts by inverse
substitution on the argument. Indeed evaluation at zero is rationally
invariant, and evaluating a translated vector at zero evaluates
the vector at its translating argument. Heisenberg covariance
transports that translating argument by the symplectic action.
Thus for the finite translate and its transported counterpart
V74M2 gives equality of the classical functions at $z$ and $dz$
up to $\tau$. The Gaussian gauge in this comparison is unchanged
because $d$ is unitary. The finite Heisenberg translations match
literally by $Y'(\pi q)=dY(q)$, with the same sign convention on
both sides. The chain rule identifies differentiation at the origin
by $P$ on the composed function with differentiation by $P'$ on
the moved function: a direction in column $b$ is sent to the sum
of directions in columns $a$ with coefficients $d_{ab}$ in
the corresponding embedding. Apply this first to a monomial in
directions, then use linearity. Thus each new summand is $\tau$
times its old partner. The weight identity
$w'(\pi q)=w(q)$ and the residue permutation finish the calculation.

For the spanning application work with the complex reading map
$L=\mathrm{atDict}\circ\mathrm{rdU}$ throughout. Starting from
the stable span by identity-cover short sections with varying
generating residues, apply the preceding formula first to the
residue move and then to the split move. Each ratio is constant
on the entire original dictionary. The resulting split readings
therefore span every $Lf$ occurring in that stable span.
The coefficient-field descent for the final class sums is K73G4.
This use requires only the target lattices realized by the original
dictionary, all of which have a target mover by its marking
condition. S74B2 also describes descent of a uniform comparison on
all complex framed points, when that stronger input is available.
\end{proof}
% END V74M3
